\section{Mengen }

Georg Kantor hat 1895 eine Menge wie folgt beschrieben:\\
Eine Menge ist eine wohldefinierte Zusammenfassung verschiedener Objekte zu einem Ganzen.\\
Beispiele:\\

\begin{tabular}[t]{lll}% l=linksbündig r=rechtsbündig
	\textbullet& $S:=\{$Berlin, Hamburg, Köln, Hannover$\}$  \\
	\textbullet &$ \mathbb{R}$:=$\{x|\;x$ ist eine reelle Zahl\}$\quad \overset{\wedge}{=}$  \glqq Alle Kommazahlen\grqq \vspace*{4mm}\\	
	& $\blacktriangleright\sqrt{2}\in\mathbb{R}, \{\sqrt{(2)}\}\subset \mathbb{R}$\vspace*{4mm}\\
	& $\blacktriangleright \emptyset:=\{ \}\subset \mathbb{R}$\vspace*{4mm}\\
	& $ \blacktriangleright\infty\notin\mathbb{R}$\vspace*{4mm}\\
	& $\blacktriangleright\{1,2,7,\mathrm{Berlin}\}\notin \mathbb{R}\vspace*{4mm}$
\end{tabular}\vspace*{2mm} \\
Wichtige Zahlenmengen:\\
\begin{tabular}[t]{lll@{\textbullet}ll}% l=linksbündig r=rechtsbündig
	&&&& $\mathbb{N}:=\{0,1,2,3,\ldots\}$ \\
	&&&& $\mathbb{Z}:=\{\ldots,-2,-1,0,1,2,\ldots\}$\\
	&&&& $\mathbb{Q}:=\Big\{\frac{p}{q}\, |\, p\in\mathbb{Z}\wedge q\in\mathbb{N}\backslash\{0\}\Big\}$ \\
	&&&& $\mathbb{C}:=\{a+\mathrm{i}b \, |\; a,b\in\mathbb{R}\; \text{und } \; \mathrm{i}^2=-1\}$	
\end{tabular}\vspace*{4mm} \\
Achtung oft werden die natürlichen Zahlen so definiert:\\
\begin{tabular}[t]{lll@{\textbullet}ll}% l=linksbündig r=rechtsbündig
	&&&& $\mathbb{N}:=\{1,2,3,\ldots\}$ \\
	&&&& $\mathbb{N}_0:=\{0,1,2,\ldots\}$\\
	
\end{tabular} \vspace*{4mm}\\
Eigenschaft:\\
$\mathbb{N}\subset\mathbb{Z}\subset\mathbb{Q}\subset\mathbb{R}\subset\mathbb{C}$\\
Mathematische Schreibweisen:\vspace*{2mm}\\
$|$ für die gilt\\
$\notin $ ist kein Element von\\
$\in$ ist Element von\\
$\{\}$ leere Menge\\



\subsection{Intervalle}

\textbf{Intervalle:}\\
seien $a,b\in\mathbb{R}$\\
$\blacktriangleright$ Abgeschlossene Intervall\\  
\begin{quote}
	$[a,b]\subseteq\mathbb{R}:=\{x\in\mathbb{R}\;|\;a\leq x \leq b\}$\vspace*{4mm}\\
	\begin{tikzpicture}
		\node at (2,1) [rotate=0,color=red] {$[$};
		\node at (2,0.5) [rotate=0, color=red] {$a$};
		\draw[->] (0,1) -- (6,1) node [below]{$\mathbb{R}$};
		\draw  (2,1) -- (4,1)[ color=red] ;
		\node at (4,1) [rotate=0,color=red] {$]$};
		\node at (4,0.5) [rotate=0, color=red] {$b$};
	\end{tikzpicture}\\
	z.B.: $[1,3]$\\
	\begin{tikzpicture}
		\node at (2,1) [rotate=0,color=red] {$[$};
		\node at (2,0.5) [rotate=0, color=red] {$1$};
		\draw[->] (0,1) -- (6,1) node [below]{$\mathbb{R}$};
		\draw  (2,1) -- (4,1)[ color=red] ;
		\node at (4,1) [rotate=0,color=red] {$]$};
		\node at (4,0.5) [rotate=0, color=red] {$3$};
	\end{tikzpicture}\\
\end{quote}
$\blacktriangleright$ Rechts halboffene Intervall\\ 
\begin{quote}
	$[a,b[\subseteq\mathbb{R}:=\{x\in\mathbb{R}\;|\;a\leq x < b\}$\vspace*{4mm}\\
	\begin{tikzpicture}
		\node at (2,1) [rotate=0,color=red] {$[$};
		\node at (2,0.5) [rotate=0, color=red] {$a$};
		\draw[->] (0,1) -- (6,1) node [below]{$\mathbb{R}$};
		\draw  (2,1) -- (4,1)[ color=red] ;
		\node at (4,1) [rotate=0,color=black] {$[$};
		\node at (4,0.5) [rotate=0, color=black] {$b$};
	\end{tikzpicture}\\
	$b\notin[a,b[$\\
	(analog: das limks soffene Intervall: $\quad  ]a,b]$) \\
\end{quote}
$\blacktriangleright$ Offene Intervall\\ 
\begin{quote}
	$]a,b[\subseteq\mathbb{R}:=\{x\in\mathbb{R}\;|\;a<x < b\}$\vspace*{4mm}\\
	\begin{tikzpicture}
		\node at (2,1) [rotate=0,color=black] {$]$};
		\node at (2,0.5) [rotate=0, color=black] {$a$};
		\draw[->] (0,1) -- (6,1) node [below]{$\mathbb{R}$};
		\draw  (2,1) -- (4,1)[ color=red] ;
		\node at (4,1) [rotate=0,color=black] {$[$};
		\node at (4,0.5) [rotate=0, color=black] {$b$};
	\end{tikzpicture}\ \\
\end{quote}
$\blacktriangleright$ Spezielle Teilintervalle\\
\begin{quote}
	Geometrisch gedeutet ist dies eine Teilgerade der reellen Zahlen.\vspace*{4mm}\\ $]-\infty,a[\subseteq\mathbb{R}:=\{x\in\mathbb{R}\;|\;x<a\}$\vspace*{4mm}\\
	\begin{tikzpicture}
		
		
		\draw[<->] (0,1) -- (6,1) node [below]{$\mathbb{R}$};
		\draw  (0,1) -- (4,1)[ color=red] ;
		\node at (4,1) [rotate=0,color=black] {$[$};
		\node at (4,0.5) [rotate=0, color=black] {$a$};
	\end{tikzpicture}\ \\
\end{quote}
Offene Intervalle werden auch statt: $]a,b[$ auch  so dargestellt: $(a,b)$.\\
Das rechts halboffene Intervall: $[a,b)$\\
Das links halboffene Intervall: $(a,b]$\\
\subsection{Durchschnitt, Vereinigung und Differenzmenge}

\textbf{Durchschnitt, Vereinigung, Differenzmenge}\\
$\blacktriangleright$ Durchschnitt\\
\begin{quote}
	
	\tikzset{
		schraffiert/.style={pattern=horizontal lines,pattern color=#1},
		schraffiert/.default=black
	}
	\begin{tikzpicture}
		\def\kreisUeins{(0.5,2) circle [radius=1.5]}
		\def\kreisUzwei{(1.5,2) circle [radius=1.5]}
		\begin{scope} % lokal halten der Wirkung von \clip
			\clip\kreisUeins;
			\draw[schraffiert=green]\kreisUzwei;
		\end{scope}
		\draw\kreisUeins;
		\draw\kreisUzwei;
		\node[anchor=east] at (-0.3,2) {$A$};
		\node[anchor=west] at (2.2,2) {$B$};
	\end{tikzpicture}\\
	
	$\qquad \quad A\cap B$\\
	Beispiel:\\
	Bilde: $[1,2] \cap [2,3]$\\
	\begin{tikzpicture}
		
		
		\draw[->] (0,1) -- (6,1) node [below]{$\mathbb{R}$};
		
		\node at (1,1) [rotate=0,color=black] {$[$};
		\node at (1,0.5) [rotate=0, color=black] {$1$};
		\node at (2,1) [rotate=0,color=black] {$[$};
		\node at (2,0.5) [rotate=0, color=black] {$2$};
		
		
	\end{tikzpicture}\ \\
	\begin{tikzpicture}
		
		
		\draw[->] (0,1) -- (6,1) node [below]{$\mathbb{R}$};
		
		\node at (2,1) [rotate=0,color=black] {$[$};
		\node at (2,0.5) [rotate=0, color=black] {$2$};
		\node at (3,1) [rotate=0,color=black] {$]$};
		\node at (3,0.5) [rotate=0, color=black] {$3$};
		
		
	\end{tikzpicture}\ \\
	$[1,2] \cap [2,3]=\{2\}$\\
\end{quote}
$\blacktriangleright$ Vereinigung\\
\begin{quote}
	
	\tikzset{
		schraffiert/.style={pattern=horizontal lines,pattern color=#1},
		schraffiert/.default=black
	}
	
	\begin{tikzpicture}
		\def\kreisUeins{(0.5,2) circle [radius=1.5]}
		\def\kreisUzwei{(1.5,2) circle [radius=1.5]}
		\begin{scope} % lokal halten der Wirkung von \clip
			%\clip\kreisUeins;
			\draw[schraffiert=green]\kreisUzwei;
			\draw[schraffiert=green]\kreisUeins;
		\end{scope}
		\draw\kreisUeins;
		\draw\kreisUzwei;
		\node[anchor=east] at (-0.3,2) {$A$};
		\node[anchor=west] at (2.2,2) {$B$};
	\end{tikzpicture}\\
	Beispiel:\\
	Bilde: $[1,2]\cup [2,3] $\\
	\begin{tikzpicture}
		
		
		\draw[->] (0,1) -- (6,1) node [below]{$\mathbb{R}$};
		
		\node at (1,1) [rotate=0,color=black] {$[$};
		\node at (1,0.5) [rotate=0, color=black] {$1$};
		\node at (2,1) [rotate=0,color=black] {$]$};
		\node at (2,0.5) [rotate=0, color=black] {$2$};
		
		
	\end{tikzpicture}\ \\
	\begin{tikzpicture}
		
		
		\draw[->] (0,1) -- (6,1) node [below]{$\mathbb{R}$};
		
		\node at (2,1) [rotate=0,color=black] {$[$};
		\node at (2,0.5) [rotate=0, color=black] {$2$};
		\node at (3,1) [rotate=0,color=black] {$]$};
		\node at (3,0.5) [rotate=0, color=black] {$3$};
		
		
	\end{tikzpicture}\ \\
	$[1,2]\cup [2,3]=[1,3]$\\
\end{quote}
\pagebreak
$\blacktriangleright$ Differenzmenge\\
\begin{quote}
	
	\tikzset{
		schraffiert/.style={pattern=horizontal lines,pattern color=#1},
		schraffiert/.default=black
	}
	\begin{tikzpicture}
		\def\kreisUeins{(0.5,2) circle [radius=1.5]}
		\def\kreisUzwei{(1.5,2) circle [radius=1.5]}
		\begin{scope} % lokal halten der Wirkung von \clip
			\clip\kreisUeins;
			\draw[schraffiert=green]\kreisUzwei;
			\draw[schraffiert=green]\kreisUeins;
		\end{scope}
		\draw\kreisUeins;
		\draw\kreisUzwei;
		\fill[white] (1.5,2) circle (1.5);
		\draw (0.5,2) circle (1.5);
		\node[anchor=east] at (-0.3,2) {$A$};
		\node[anchor=west] at (2.2,2) {$B$};
	\end{tikzpicture}\\
	Beispiel:\\
	$[1,3] \; \backslash \;] \frac{3}{2},\frac{5}{2} ]$	\\
	\begin{tikzpicture}
		
		
		\draw[->] (0,1) -- (6,1) node [below]{$\mathbb{R}$};
		
		\node at (1,1) [rotate=0,color=black] {$[$};
		\node at (1,0.5) [rotate=0, color=black] {$1$};
		\node at (3,1) [rotate=0,color=black] {$]$};
		\node at (3,0.5) [rotate=0, color=black] {$3$};
		
		
	\end{tikzpicture}\ \\
	\begin{tikzpicture}
		
		
		\draw[->] (0,1) -- (6,1) node [below]{$\mathbb{R}$};
		
		\node at (1.5,1) [rotate=0,color=black] {$]$};
		\node at (1.5,0.5) [rotate=0, color=black] {$\frac{3}{2}$};
		\node at (2.5,1) [rotate=0,color=black] {$]$};
		\node at (2.5,0.5) [rotate=0, color=black] {$\frac{5}{2}$};
		
		
	\end{tikzpicture}\ \\
	$[1,3] \;\backslash \;] \frac{3}{2},\frac{5}{2} ]=[1,\frac{3}{2}] \cup ]\frac{5}{2}, 3]$	\\
	
	\begin{tikzpicture}
		
		
		\draw[->] (0,1) -- (6,1) node [below]{$\mathbb{R}$};
		
		\node at (1,1) [rotate=0,color=black] {$[$};
		\node at (1,0.5) [rotate=0, color=black] {$1$};
		\node at (1.5,1) [rotate=0,color=black] {$]$};
		\node at (1.5,0.5) [rotate=0, color=black] {$\frac{3}{2}$};
		\node at (2.5,1) [rotate=0,color=black] {$]$};
		\node at (2.5,0.5) [rotate=0, color=black] {$\frac{5}{2}$};
		\node at (3,1) [rotate=0,color=black] {$]$};
		\node at (3,0.5) [rotate=0, color=black] {$3$};
	\end{tikzpicture}\ \\	
\end{quote}
\subsection{Aufgabe zu Mengen 1}

\textbf{Aufgabe:}\\
Bestimme die folgende Menge $\mathbb{M}$:\\
$\mathbb{M}=\{x\in \mathbb{Z}\;|\; \text{$ x $  ist durch 3 teilbar} \}\cap \{x\in \mathbb{Z}\ \;|\;\text{$ x $ ist durch 4 teilbar } \}$\vspace*{2mm}\\
Für die Menge der Zahlen die durch 3 teilbar sind gilt:\\
\begin{align*}
	\{3x\;|\;x\in \mathbb{Z}\}\\
\end{align*}
Für die Menge der Zahlen die durch 4 teilbar sind gilt:\\
\begin{align*}
	\{4x\;|\;x\in \mathbb{Z}\}\\
\end{align*}
Für welche Menge gilt dann, dass $x$ durch 3 und durch 4 teilbar ist?\\
Das Produkt von 4 und 3 ist immer durch 3 und 4 teilbar. Also gilt:\\
$\mathbb{M}=\{x\in\mathbb{Z}\; | \; x\;\text{ist durch 3 und 4 teilbar} \}=\{12x\;|\;x\in\mathbb{Z}\}$\vspace*{2mm}\\




\subsection{Aufgabe zu Mengen 2}

\textbf{Aufgabe:}

Bestimme die folgende Menge $\mathbb{M}$:

\begin{align*}
	\mathbb{M}
	= \bigcup_{m\in\mathbb{N}\setminus\{0\}}
	\{x\in\mathbb{Q}\mid m\cdot x\in\mathbb{Z}\}
\end{align*}

\textbf{Lösung:}

Betrachten wir zunächst den Fall $m=2$:

\begin{align*}
	M_2=\{x\in\mathbb{Q}\mid 2x\in\mathbb{Z}\}
\end{align*}

Welche rationalen Zahlen ergeben mit 2 multipliziert
eine ganze Zahl?

Beispiele:

\begin{align*}
	\frac{1}{2}\cdot 2 &= 1\\
	-\frac{1}{2}\cdot 2 &= -1\\
	\frac{7}{2}\cdot 2 &= 7
\end{align*}

Damit erhalten wir:

\begin{align*}
	M_2
	= \left\{\frac{n}{2}\mid n\in\mathbb{Z}\right\}
\end{align*}

Für ein beliebiges festes
$m\in\mathbb{N}\setminus\{0\}$ gilt entsprechend:

\begin{align*}
	M_m
	&= \{x\in\mathbb{Q}\mid mx\in\mathbb{Z}\}\\
	&= \left\{\frac{n}{m}\mid n\in\mathbb{Z}\right\}
\end{align*}

Vereinigt man diese Mengen über alle positiven
natürlichen Zahlen $m$, so erhält man sämtliche
rationalen Zahlen.

Denn jede rationale Zahl lässt sich als Bruch
$\frac{n}{m}$ mit $n\in\mathbb{Z}$ und
$m\in\mathbb{N}\setminus\{0\}$ darstellen.

Somit gilt:

\begin{align*}
	\boxed{
		\mathbb{M}
		= \bigcup_{m=1}^{\infty} M_m
		= \bigcup_{m=1}^{\infty}
		\left\{\frac{n}{m}\mid n\in\mathbb{Z}\right\}
		= \mathbb{Q}
	}
\end{align*}



\subsection{Aufgabe zu Mengen 3}

Bestimme die folgende Menge:\\ 
\begin{align*}
	\mathbb{M}:\;\overset{\bigcap}{\scriptscriptstyle  m \in \mathbb{R}}\{x\in\mathbb{R}\;|\;(x-m)(x-2)(x-3)=0  \}
\end{align*}
Für $m=1$ gilt:\\
\begin{align*}
	m=1:\{x\in\mathbb{R}\;|\;(x-1) (x-2) (x-3)=0 \}=\{1,2,3\}
\end{align*}
Für $m=\sqrt{2}$ gilt:\\
\begin{align*}
	m=\sqrt{2}:\{x\in\mathbb{R}\;|\;(x-\sqrt{2})(x-2)(x-3) \}=\{\sqrt{2},2,3\}
\end{align*}
Für m allgemein gilt:\\
\begin{align*}
	m:\{x\in\mathbb{R}\;|\;(x-m)(x-2)(x-3) \}=\{m,2,3\}
\end{align*}
Die Schnittmenge:\\
\begin{align*}
	\mathbb{M}:\;\overset{\bigcap}{\scriptscriptstyle  m \in \mathbb{R}}\{x\in\mathbb{R}\;|\;(x-m)(x-2)(x-3)=0  \}=\{2,3\}
\end{align*}
(weil für x=2 und x=3 und für m=2 und m=3 eine Lösung existiert)

